\chapter{Conclusion and Future Work}
\label{chap:conclusion}

This thesis investigated the problem of connecting wildfire-risk prediction with wildfire-loss assessment. The central motivation was that much of the recent wildfire literature remains focused on hazard prediction, including ignition likelihood, susceptibility, spread, detection, and fire-danger estimation, while fewer studies translate predictive outputs into economic or ecological consequence estimates. To address this gap, the thesis developed a county-level wildfire-loss modeling framework using publicly accessible datasets.

The literature review showed that recent research has made substantial progress in wildfire prediction, remote sensing, deep learning, ecological vulnerability assessment, and AI-supported decision systems. At the same time, it highlighted that integrated wildfire-loss modeling remains comparatively underdeveloped. This is especially true for scalable frameworks that connect hazard-relevant environmental variables with monetary loss-oriented targets and geographic validation.

The main methodological contribution of the thesis is a complete FEMA/NRI-based county-level wildfire risk prediction and assessment pipeline. The pipeline begins with a county-level wildfire Expected Annual Loss target in monetary units. It then adds county-level exposure, value, vulnerability, and resilience predictors, followed by environmental predictors derived from topography, land cover, and warm-season climate summaries. The resulting feature table is modeled using regression algorithms and evaluated under both random and state-wise grouped validation.

The experimental results show that the county-only baseline provides a meaningful but limited predictive signal. Environmental augmentation produces a substantial improvement, indicating that wildfire losses are strongly associated with terrain, land cover, and climate conditions in addition to county-level exposure and resilience variables. The strongest performance was obtained with the combined Random Forest model. State-wise grouped validation further showed that although geographic generalization is more difficult than random splitting, the combined model remains useful when evaluated across held-out states.

The final pipeline output demonstrates that the model can be used for county-level wildfire risk assessment, not only for regression benchmarking. By converting predicted log-transformed wildfire EAL back into monetary units, the workflow produces predicted wildfire-loss estimates, county rankings, and percentile-based risk categories. The thesis also adds an exploratory ecological risk-screening layer that combines predicted wildfire-risk percentile with natural vegetation exposure. This makes the thesis contribution consequence-oriented: the output is an interpretable wildfire-risk assessment product in monetary units, complemented by an ecological exposure-screening perspective, rather than only a wildfire occurrence or hazard prediction.

The thesis also documented an earlier Creek Fire tensor-based workflow as a supporting baseline and future extension. This workflow demonstrates fine-scale data engineering using fire detections, a fire perimeter, static geospatial predictors, and time-varying climate features. However, the main defendable M.S. contribution is the FEMA/NRI county-level wildfire risk/loss assessment pipeline, because it is directly connected to monetary wildfire-loss estimation.

Several limitations remain. The FEMA NRI wildfire EAL target is spatially aggregated and does not capture event-level or asset-level damage mechanisms directly. The environmental augmentation currently focuses on the contiguous United States and does not fully address all county-equivalent units. The validation strategy includes state-wise grouped validation, but additional spatial validation designs and uncertainty analysis would strengthen the model further. In addition, future work should incorporate more detailed explainability methods, richer loss data, and ecological consequence indicators.

Future research can extend this work in three main directions. First, the modeling framework can be improved with additional wildfire-relevant predictors, uncertainty quantification, regional calibration, and more advanced explainability methods. Second, the data foundation can be enriched using event-level wildfire losses, building-level exposure, infrastructure networks, suppression-cost data, smoke-health impacts, and ecological vulnerability indicators. Third, the exploratory Creek Fire workflow can be developed into a fine-scale spatio-temporal modeling pipeline that complements the county-level FEMA/NRI framework.

The empirical contribution of the thesis remains primarily economic: it models county-level wildfire Expected Annual Loss in monetary units. Ecological losses are addressed through the literature review, through the inclusion of environmental predictors, and through an exploratory ecological proxy based on natural vegetation exposure. However, ecological losses are not directly estimated as a separate damage target. Future work should extend the current framework by integrating ecological vulnerability, burn severity, vegetation recovery, biodiversity exposure, and ecosystem-service loss indicators.

The supporting Vietnam study further demonstrates that fine-scale wildfire prediction can benefit from dynamic fire-growth features, threshold-aware evaluation, and SHAP-based interpretability, while the FEMA/NRI pipeline extends the thesis toward county-level monetary loss assessment.

Part of this work was successfully presented at DCHPC 2026 through the paper entitled ``Mapping Forest Fire Risks Using Deep Learning: A Case Study in An Giang Province, Vietnam.'' This conference presentation supports the broader research direction of the thesis by demonstrating fine-scale wildfire prediction using deep learning, while the main thesis contribution extends the work toward county-level monetary wildfire-loss assessment. The corresponding certificate is included in Appendix~\ref{app:conference_certificate}, and the associated paper is expected to appear as a Web of Science/Scopus paper.

The main novelty of the thesis lies in the general workflow and insights rather than in the specific territory used for implementation. Unlike previous wildfire-prediction studies that often stop at hazard, susceptibility, or burned-area prediction, this thesis connects wildfire-relevant environmental conditions with exposure, vulnerability, resilience, and monetary expected loss. The new methodological contribution is an integrated wildfire-loss assessment pipeline that combines environmental augmentation, machine-learning regression, geographic validation, county-level risk ranking, and exploratory ecological risk screening.

The key insights gained are threefold. First, environmental predictors such as precipitation, vapor pressure deficit, land-cover composition, elevation, slope, and wind contain strong information for wildfire-loss modeling. Second, geographically stricter validation gives a more conservative and realistic estimate of model transferability than random splitting alone. Third, monetary wildfire-loss assessment can be complemented by ecological exposure screening, which helps identify counties where high predicted wildfire risk overlaps with high natural vegetation exposure. These insights make the framework transferable to other regions where comparable wildfire hazard, exposure, vulnerability, environmental, and loss-related datasets are available.

Overall, the thesis demonstrates that publicly accessible data can be used to construct a complete county-level wildfire risk prediction and assessment pipeline focused on monetary expected loss. The results support the broader thesis direction of moving from hazard-oriented wildfire prediction toward integrated, consequence-oriented wildfire-loss assessment.